\section{Finite stability when no convergence modulus is available}
\label{sec:stability}

An explicit modulus of convergence is excellent input when it is available.
It should not, however, be made a prerequisite for every use of a bounded
sequence. Many arguments inspect only a finite interval of indices. For
such arguments a weaker, precisely quantified conclusion suffices. This
section gives one arithmetic selection principle that handles several
analytic observations and several finite exceptional sets at the same
time. The principle is a finite version of standard proof-mining ideas
\cite{ProofMining}; the additional purpose here is to expose exactly the
finite hypotheses and the quantifier-free invariants needed for iPRA.

\subsection{The finite data and the common-window theorem}

Fix a primitive recursive function $g:\NN\to\NN$ and an initial index $b$.
No monotonicity assumption on $g$ is required. Define
\begin{equation}\label{eq:endpoints}
 n_0=b,\qquad n_{i+1}=n_i+g(n_i)+1,
 \qquad W_i=[n_i,n_i+g(n_i)]\cap\NN.
\end{equation}
The windows are disjoint and ordered. All the endpoints are primitive
recursive, for a fixed primitive recursive definition of $g$.

For $j<m$, let $a_j(t)=a(j,t)$ for one fixed primitive recursive rational
evaluator $a$, let
$\epsilon_j>0$ be rational, and let $V_j\geq0$ be rational. For $\ell<r$,
let $E_\ell(t)=E(\ell,t)$ for one fixed primitive recursive Boolean
evaluator $E$, with natural
budget $C_\ell$. A successful window is one on which
\begin{align}
 |a_j(p)-a_j(q)|&\leq\epsilon_j
   &&(j<m,\ p,q\in W_i),\label{eq:stable-window}\\
 E_\ell(t)&=0 &&(\ell<r,\ t\in W_i).\label{eq:empty-window}
\end{align}
Additional natural parameters of these fixed evaluators are permitted;
the indices $j,\ell$ do not range over arbitrary program codes.
Every test in these formulas is finite and decidable. Use rational
cross-multiplication for the comparisons, not floating-point arithmetic.
Put
\begin{equation}\label{eq:stability-budget}
 K_j=\left\lfloor\frac{V_j}{\epsilon_j}\right\rfloor,
 \qquad K=\sum_{j<m}K_j+\sum_{\ell<r} C_\ell,
 \qquad T=n_{K+1}.
\end{equation}

\begin{theorem}[A common window from additive finite budgets]
\label{thm:common-window}
Suppose the following \emph{finite} inequalities hold:
\begin{align}
 \sum_{t<T}|a_j(t+1)-a_j(t)|&\leq V_j &&(j<m),\label{eq:finite-variation}\\
 \sum_{t<T} E_\ell(t)&\leq C_\ell &&(\ell<r).
 \label{eq:finite-exceptions}
\end{align}
Then a bounded search among $i=0,\ldots,K$ finds a successful window.
In particular its starting point is at most $n_K$. For fixed primitive
recursive evaluators $a_j,E_\ell,g$, this is an iPRA theorem with
quantifier-free finite input and output predicates.
\end{theorem}

\begin{proof}
For a fixed component $j$, call a window $j$-bad if the first test fails
for that component. Bounded search then supplies $p,q\in W_i$ with
$|a_j(p)-a_j(q)|>\epsilon_j$. Swapping $p,q$ if necessary and telescoping
gives
\[
 \epsilon_j<|a_j(p)-a_j(q)|
 \leq\sum_{t=p}^{q-1}|a_j(t+1)-a_j(t)|
 \leq\sum_{t=n_i}^{n_i+g(n_i)-1}|a_j(t+1)-a_j(t)|.
\]
For different $i$ these ranges of increments are disjoint. Consequently,
if $L$ windows are $j$-bad, then
$L\epsilon_j<\sum_{t<T}|a_j(t+1)-a_j(t)|\leq V_j$ when $L>0$.
There can therefore be at most $K_j$ such windows. This bound is slightly
generous when $V_j/\epsilon_j$ is an integer; generosity avoids a special
case and does no harm.

For a fixed exceptional predicate $E_\ell$, every window failing its
emptiness test contains a least exceptional index. The chosen indices
are distinct because the windows are disjoint, so there are at most
$C_\ell$ such windows by \eqref{eq:finite-exceptions}.

If all $K+1$ windows failed, assign each window to its least failed test,
in a fixed order of the $m+r$ tests. Each assigned class has size at most
its budget by the preceding two paragraphs. Summing the class counts
would give $K+1\leq K$, a contradiction. Because the tests are decidable,
this proves a bounded positive existence statement: scan the $K+1$
windows and retain the first success. If the scan returned failure,
the finite argument just given would contradict its recorded failures.
No double-negation elimination for an unbounded proposition is used.

All proof loops have visible finite bounds. The triangle inequality is
proved by induction on the number of increments; disjoint-range sums are
proved by induction on the number of selected windows; the last counting
identity is proved by induction on the list of assigned tests. Each
invariant is a rational inequality, a bounded conjunction, or an equality
of finite sums. After rational coding and Booleanization these are
quantifier-free formulas. The endpoint $T$ is fixed before these
inductions. Thus the proof has the required iPRA form.
\end{proof}

The use of variation through $T$ intentionally overestimates what is
needed: the last tested window ends at $T-1$. This harmless extra
increment makes the interface easy to compose. For $m=r=0$, $K=0$ and
the first window succeeds vacuously. A zero-length window also behaves
correctly: its oscillation is zero, while its exceptional-index test
still inspects its single index.

\subsection{Monotone sequences, finite variation, and simultaneous use}

If $a(t)$ is nondecreasing and $0\leq a(t)\leq M$, then for every finite
$T$ its variation telescopes to $a(T)-a(0)\leq M$. With
$\epsilon=2^{-k}$, Theorem~\ref{thm:common-window} therefore gives
\[
 \exists i\leq \lfloor M2^k\rfloor\quad
 \forall p,q\in[n_i,n_i+g(n_i)]\quad |a(p)-a(q)|\leq2^{-k}.
\]
The displayed quantifiers over $i,p,q$ are bounded; the search returns
the particular $i$. This is often called \emph{metastability}. Its
meaning is finite-window stability, not eventual stability on an
infinite tail.

The same theorem covers nonmonotone sequences with a proved finite
variation bound. It also synchronizes several sequences without first
finding a stable interval for one sequence and then searching inside it
for another. Their budgets add. Exceptional sets can be included in
the same selection. This matters when a proof needs, at one common
scale, a small oscillation in several approximants and avoidance of a
finite family of bad indices.

For example, suppose two monotone rational sequences take values in
$[0,1]$, both tolerances are $1/4$, and a decidable exceptional predicate
has at most two values in the inspected prefix. Then $K=4+4+2=10$.
Eleven candidate windows suffice for a common successful one. If
$g(n)=n+1$ and $b=0$, then $n_i=2^{i+1}-2$, so the returned start is
at most $2046$. This is a guaranteed bound; a particular input usually
succeeds much earlier.

\paragraph{Relation to the existing Roth packet.}
The repository already observes that a bound on the number of good
approximants yields a finite window with no such approximants
\cite{ProveIt}. Theorem~\ref{thm:common-window} does not claim that
observation as new, nor does it prove the missing quantitative Roth
estimate. It generalizes the common last step to multiple exceptional
predicates and analytic variation budgets, and supplies one precise
finite interface usable by other proofs. An absent cardinality estimate
remains an absent hypothesis.

\subsection{Why the output cannot be silently upgraded to a modulus}

The difference between a stable window and a Cauchy modulus is essential.
Let $H(e,t)$ be the decidable predicate that a fixed universal machine
with input $e$ has halted within $t$ steps. Bounded simulation is primitive
recursive. The sequence $a_e(t)=H(e,t)$ is nondecreasing and lies in
$[0,1]$. If a computable procedure uniformly returned a Cauchy cutoff
$M(e)$ for tolerance $1/2$, simulating through $M(e)$ would decide
halting: if the machine had not halted then, it could never halt later,
since that would change the sequence from $0$ to $1$ after the cutoff.
Thus even this very simple uniform family has no uniform computable
Cauchy modulus, although the finite-window theorem applies to it.

One can also see the obstruction in a single fixed primitive recursive
rational sequence. Define
\[
 s_t=\sum_{e\leq t}4^{-(e+1)}H(e,t).
\]
It is nondecreasing, primitive recursive, and bounded by $1/3$. Suppose
it had a computable Cauchy modulus. In the ordinary external real-number
interpretation, that modulus would compute its limit $s$. To recover the
halting status of machine $e$, first subtract the contributions of
machines $0,\ldots,e-1$, which have already been recovered. Put
$w=4^{-(e+1)}$. If $e$ never halts, the remaining value is at most $w/3$;
if it halts, the remaining value is at least $w$. An approximation with
error at most $w/6$ separates these cases, for example at threshold
$2w/3$. Recursing through $e$ decides the halting problem, a contradiction.
This establishes that a global modulus cannot in general be extracted
even for a single primitive recursive bounded monotone rational sequence.

The external impossibility argument is not used as an axiom in any
positive iPRA proof. It diagnoses a proposed change of output contract.
Nor does it license a blanket claim that an arbitrary mathematical
theorem is unprovable in iPRA. One must first identify its exact input
representation and exact conclusion.

\subsection{The Lean and arithmetic boundaries of the selector}

The shipped \texttt{FiniteCore.lean} defines the bounded scan
\texttt{firstHit}, disjoint endpoints, and finite contracts. It proves
that a successful index within the budget makes the scan succeed, that
the scan stays within the budget, and that bounded contracts compose.
It also proves a natural-valued energy termination lemma, a useful
alternative implementation when the failed-test budget is already an
integer potential.

These are compiled ambient Lean theorems. The rational variation and
multi-component counting proof above has a formalization outline, not a
completed Lean proof. To finish it, represent the rational arrays by
finite lists; prove interval telescoping, the disjoint-sum inequality,
and the finite assignment count; and then apply
\texttt{firstHit\_complete}. To finish the \emph{object-theory} proof,
replace the arrays by natural codes and construct the corresponding
\texttt{Prf PRAAx} derivations using the bounded invariants just listed.
Compilation of the ambient selector alone does not perform that last
step.
